\documentclass[11pt]{letter}

\usepackage[a4paper, margin=1in]{geometry}
\usepackage[T1]{fontenc}
\usepackage{lmodern}
\usepackage{hyperref}
\hypersetup{colorlinks=true, urlcolor=black, linkcolor=black}

\setlength{\parindent}{0pt}
\setlength{\parskip}{10pt}

\begin{document}

\begin{center}
{\Large \textbf{Ivan Evdokimov, Ph.D.}}\\[4pt]
{\small Manchester, UK \; $|$ \; ivevdm@gmail.com \; $|$ \; \href{https://linkedin.com/in/ivan-evdokimov}{linkedin.com/in/ivan-evdokimov} \; $|$ \; \href{https://github.com/ivan020}{github.com/ivan020}}
\end{center}

\vspace{10pt}

Dear Hiring Manager,

I am applying for this role because the University of Exeter is one of the UK’s leading universities for Data and Computer Science. This position offers a rare chance to combine rigorous research with production-grade data infrastructure, which is exactly what I have built my career on. Data management principles are essential today, as data has become a critical commodity that drives innovation across many fields.
I hold a PhD in Computational Finance from the University of Essex. My doctoral work developed a PyTorch-based transfer learning framework for Bayesian model averaging in financial forecasting. This resulted in a first-authored peer-reviewed journal paper, three co-authored conference papers, and an open-source package that has been adopted beyond my own research group, proving the work was both reproducible and useful to others.

Python is my primary programming language, and I have used it extensively for backend services, machine learning models, and LLM integrations. I am also fluent in C++, having sole-authored a modular agent-based simulation during my PhD and later migrating statistical algorithms from R to C++ to achieve an 11x performance improvement. I have additionally picked up Java, Golang, and TypeScript as different projects required them. I build CI/CD pipelines with automated testing and code review, and I use Docker and infrastructure-as-code practices to ensure my work is reliable and reproducible. When needed, I quickly learn new tools and technologies, such as LangGraph for agent workflows or Hugging Face and vLLM for fine-tuning and deploying LLaMA models.
At UK Data Service, I trained and deployed PyTorch and scikit-learn models for automated metadata classification and disclosure risk assessment, achieving over 95\% F1 accuracy on sensitive social science datasets. When existing R-based algorithms could not handle growing data volumes, I researched and implemented a C++ replacement that enabled near real-time processing of datasets exceeding 500,000 records. I have delivered lectures and labs in Data Structures, C/C++, and Machine Learning to over 100 students, achieving an average feedback score of 4.8 out of 5. This experience, along with my publication record, shows my ability to present technical work to both academic and non-technical audiences.

I have worked as part of research teams on multi-university collaborations and handled confidential and sensitive data, supporting responsible curation and controlled access to public research data. My work distributing retraining pipelines across AWS Lambda and Step Functions has required thinking about parallelism and computational efficiency at scale. I have communicated technical concepts to peer researchers, students, and non-technical end users, including building a mobile application with OpenAI API integration for field technicians, which saved an estimated three weeks per year of manual data collection.
In my current role, I manage full-stack development, MLOps infrastructure, and CI/CD across a multi-team platform with weekly releases. I balance pipeline development with administrative demands and competing deadlines by treating documentation, monitoring, and code review as part of the core work. My open-source contributions and conference co-authorship show my engagement with the wider research community. Having taught and worked with colleagues from diverse backgrounds, I understand the importance of an inclusive culture in enabling everyone to participate fully in research.

Kind regards,\\
Ivan Evdokimov

\end{document}

